% EVAL FIXTURE — deliberately flawed. Invented content, not a template.
% Failure under test: stage 4, Step-1 sufficiency. Stages 1-3 pass: the question is
% nameable, the opening carries an instability, vision and aim are separate, the three
% objectives are knowledge claims. What is missing from B1 is everything that makes the
% approach sound convincing and the PI sound inevitable -- it is all in b2.tex, which
% no Step-1 panel member ever opens.
\documentclass[11pt,a4paper]{article}
\usepackage[b1]{ercformatting}
\usepackage{ercreview}

\title{Why Cold Adaptation Reverses in Small Populations}
\acronym{REVERSAL}
\author[Vey]{M. Vey}
\institution{Institute of Evolutionary Biology, Example University, Country}

\begin{document}
\maketitle

\begin{abstract}
Cold-adapted alpine plant populations are expected to become more cold-tolerant as
they retreat upslope. Two long-term census series show the opposite: the smallest,
highest populations are losing cold tolerance. Either the selection gradient reverses
at small population size, or drift overwhelms it. This project distinguishes the two.
\end{abstract}

\startsectiona

\section{A prediction that fails at the range edge}
Adaptation to cold should strengthen where cold is strongest. Twenty years of census
data from two independent alpine transects show the reverse: in populations below
roughly 200 individuals, measured cold tolerance declines even as mean temperature
falls. Both series cannot be artefacts of the same sampling design, because the two
transects were designed and run independently.

The field currently has no account of this. The standard model predicts strengthening
adaptation and is silent about population size; the drift-load literature predicts
weakened adaptation but says nothing about why the threshold sits near 200. Whichever
explanation holds, the conservation inference drawn from these censuses for the past
decade is wrong in one direction or the other.

Why it has stayed unresolved is a matter of measurement, not of interest. Separating
reversed selection from drift requires knowing the selection gradient and the effective
population size in the same population in the same season, and until recently the
second quantity could not be estimated at this sample size.

\section{Vision and aim}
\textbf{Vision.} A theory of adaptation that holds at the range edge, where populations
are small, rather than only in the large populations the theory was built on.

\textbf{Aim.} Within five years, to determine whether loss of cold tolerance in small
alpine populations is caused by a reversal of the selection gradient or by drift.

\section{Objectives}
\textbf{O1.} Determine whether the selection gradient on cold tolerance reverses sign
below a population-size threshold.

\textbf{O2.} Identify whether the threshold is set by effective population size or by
census size, which the standard model conflates.

\textbf{O3.} Test whether the reversal, where it occurs, is predictable from genomic
load alone.

\section{State of the art}
Selection-gradient estimation and drift-load estimation have developed separately. No
study has applied both to the same population.

\section{Approach}
The three objectives are addressed by a common design across a population-size
gradient. Details of the methodology are in Part II.

\section{Why now, why me}
The methodological barrier named above was removed in 2024. My group's work is
described in Part II, Section~b.

\section{Impact}
If selection reverses, the conservation prioritisation of small high-altitude
populations inverts, because the mechanism -- and not the taxon -- is what transfers to
other range-edge systems.

\startsectionb

\cvsection*{Personal Details}
ORCID: 0000-0003-0000-0000

\cvsubsection{Current position(s)}
2019 -- present, Group Leader, Example University

\cvsubsection{Ten selected outputs}
1. Vey M. et al. (2023). Cold tolerance along an alpine transect. \emph{Example J.
Ecol.}
2. Vey M., Ora K. (2022). Effective size estimation in small plant populations.
\emph{Example Genet.}

\end{document}
